\documentclass[table]{beamer}
%[]中可以使用draft、handout、screen、transparency、trancompress、compress等参数

%指定beamer的模式与主题
\mode<presentation>
{
  \usetheme{Madrid}
%\usetheme{Boadilla}
%\usecolortheme{default}
%\usecolortheme{orchid}
%\usecolortheme{whale}
%\usefonttheme{professionalfonts}
}

%\usetheme{Madrid}
%这里还可以选择别的主题：Bergen, Boadilla, Madrid, AnnArbor, CambridgeUS, Pittsburgh, Rochester, Warsaw, ...
%有导航栏的Antibes, JuanLesPins, Montpellier, ...
%有内容的Berkeley, PaloAlto, Goettingen, Marburg, Hannover, ...
%有最小导航栏的Berlin, Ilmenau, Dresden, Darmstadt, Frankfurt, Singapore, Szeged, ...
%有章和节表单的Copenhagen, Luebeck, Malmoe, Warsaw, ...

%\usecolortheme{default}
%设置内部颜色主题(这些主题一般改变block里的颜色);这个主题一般选择动物来命名
%这里还可以选择别的颜色主题,如默认的和有特别目的的颜色主题default,structure,sidebartab,全颜色主题albatross,beetle,crane,dove,fly,seagull,wolverine,beaver

%\usecolortheme{orchid}
%设置外部颜色主题(这些主题一般改变title里的颜色);这个主题一般选择植物来命名
%这里还可以选择别的颜色主题,如默认的和有特别目的的颜色主题lily,orchid,rose

%\usecolortheme{whale}
%设置字体主题;这个主题一般选择海洋动物来命名
%这里还可以选择别的颜色主题,如默认的和有特别目的的颜色主题whale,seahorse,dolphin

%\usefonttheme{professionalfonts}
%类似的还可以定义structurebold,structuresmallcapsserif,professionalfonts

% 控制 beamer 的风格,可以根据自己的爱好修改
%\usepackage{beamerthemesplit} %使用 split 风格
%\usepackage{beamerthemeshadow} %使用 shadow 风格
%\usepackage[width=2cm,dark,tab]{beamerthemesidebar}

%插入音标
%\usepackage{tipa}
%\AtBeginDocument{
  %\renewcommand\textipa{\fontencoding{T3}\selectfont}
%}
%\AtBeginDocument{
  %\renewcommand\textipa[2][r]{{\fontfamily{cm#1}\tipaencoding #2}}
%}
%\renewenvironment{IPA}[1][r]
 %{\fontfamily{cm#1}\tipaencoding}
 %{}

% 设定英文字体
%\usepackage{fontspec}
% Fix bugs for fontspec in TeXLive2015
\ifdefined\suppressfontnotfounderror
  \expandafter\let\csname xetex_suppressfontnotfounderror:D\endcsname
    \suppressfontnotfounderror
\else
  \expandafter\let\csname xetex_suppressfontnotfounderror:D\endcsname
    \luatexsuppressfontnotfounderror
\fi
\usepackage[no-math]{fontspec}
\setmainfont{Times New Roman}
\setsansfont{Arial}
\setmonofont{Courier New}

% 设定中文字体
\usepackage[BoldFont,SlantFont,CJKchecksingle,CJKnumber]{xeCJK}
%\setCJKmainfont[BoldFont={Adobe Heiti Std},ItalicFont={Adobe Kaiti Std}]{Adobe Song Std}
\setCJKmainfont[BoldFont={Adobe Heiti Std},ItalicFont={Adobe Kaiti Std}]{WenQuanYi Micro Hei}
\setCJKsansfont{Adobe Heiti Std}
\setCJKmonofont{Adobe Fangsong Std}
\punctstyle{hangmobanjiao}

\defaultfontfeatures{Mapping=tex-text}
\usepackage{xunicode}
\usepackage{xltxtra}

\XeTeXlinebreaklocale "zh"
\XeTeXlinebreakskip = 0pt plus 1pt minus 0.1pt

\usepackage{setspace}
\usepackage{colortbl,xcolor}
\usepackage{hyperref}
%\hypersetup{xetex,bookmarksnumbered=true,bookmarksopen=true,pdfborder=1,breaklinks,colorlinks,linkcolor=blue,filecolor=black,urlcolor=cyan,citecolor=green}
\hypersetup{xetex,bookmarksnumbered=true,bookmarksopen=true,pdfborder=1,breaklinks,colorlinks,linkcolor=cyan,filecolor=black,urlcolor=blue,citecolor=green}

% 插入图片
\usepackage{graphicx}
\graphicspath{{figures/}}
% 图文混排
%\usepackage{picins}
\usepackage{floatflt}

% 可能用到的包
\usepackage{amsmath,amssymb}
%插入多媒体
%\usepackage{media9}
%\usepackage{movie15}
\usepackage{multimedia}
\usepackage{multicol}
\usepackage{multirow}

% 定义一些自选的模板，包括背景、图标、导航条和页脚等，修改要慎重
% 设置背景渐变由10%的红变成10%的结构颜色
%\beamertemplateshadingbackground{red!10}{structure!10}
%\beamertemplatesolidbackgroundcolor{white!90!blue}
% 使所有隐藏的文本完全透明、动态，而且动态的范围很小
\beamertemplatetransparentcovereddynamic
% 使itemize环境中变成小球，这是一种视觉效果
\beamertemplateballitem
% 为所有已编号的部分设置一个章节目录，并且编号显示成小球
\beamertemplatenumberedballsectiontoc
% 将每一页的要素的要素名设成加粗字体
\beamertemplateboldpartpage

% item逐步显示时,使已经出现的item、正在显示的item、将要出现的item呈现不同颜色
\def\hilite<#1>{
 \temporal<#1>{\color{gray}}{\color{blue}}
    {\color{blue!25}}
}

\renewcommand{\today}{\number\year 年 \number\month 月 \number\day 日}

%五角星
\usepackage{MnSymbol}

%去除图表标题中的figure等
\usepackage{caption}
\captionsetup{labelformat=empty,labelsep=none}

\usepackage{tabu}
\usepackage{multirow}
%表格自动换行
\usepackage{tabularx} 

% 千分号
%\usepackage{textcomp}

%罗马数字
\makeatletter
\newcommand{\rmnum}[1]{\romannumeral #1}
\newcommand{\Rmnum}[1]{\expandafter\@slowromancap\romannumeral #1@}
\makeatother

%分栏
\usepackage{multicol}

%\usepackage{enumitem}
%\usepackage{enumerate}

%键盘
\usepackage{keystroke}

%心形
%\usepackage{fdsymbol}

%插入源代码
\usepackage{listings}
\lstset{
  language=perl,                  % 程序语言名称：TeX, Perl, R, sh, bash, Awk
  basicstyle=\normalsize\tt,      %\tt指monospace字体族，程序源代码使用此族字体表示更加美观
  numbers=left,                   % 行号位置(左侧)
  numberstyle=\small,             % 行号字体的字号
  stepnumber=1,                   % 行号的显示步长
  numbersep=5pt,                  % 行号与代码间距
  backgroundcolor=\color{white},  % 背景色;需要 \usepackage{color}
  showspaces=false,               % 不显示空格
  showstringspaces=false,         % 不显示代码字符串中的空格标记
  showtabs=false,                 % 不显示 TAB
  tabsize=4, 
  frame=shadowbox,                % 把代码用带有阴影的框圈起来
  captionpos=b,                   % 标题位置
  breaklines=true,                % 对过长的代码自动断行
  breakatwhitespace=false,        % 断行只在空格处
  extendedchars=false,            % 解决代码跨页时，章节标题，页眉等汉字不显示的问题
  %escapeinside={\%*}{*},         % 跳脱字符，添加注释，暂时离开 listings 
  %escapeinside=``,
  commentstyle=\color{red!50!green!50!blue!50}\tt,  %浅灰色的注释
  rulesepcolor=\color{red!20!green!20!blue!20},     %代码块边框为淡青色
  keywordstyle=\color{blue!70}\bfseries\tt,         %代码关键字的颜色为蓝色，粗体
  identifierstyle=\tt,
  stringstyle=\tt,                % 代码字符串的特殊格式
  keepspaces=true,
  breakindent=1em,
  %breakindent=22pt,
  %breakindent=4em,
  breakautoindent=true,
  flexiblecolumns=true,
  aboveskip=1em,                  %代码块边框
  xleftmargin=2em,
  xrightmargin=2em
}

%\setbeamercolor{alerted text}{fg=magenta}
\setbeamercolor{bgcolor}{fg=yellow,bg=cyan}
%\setbeamercolor{itemize/enumerate body}{fg=green}

\begin{document}

%\includeonlyframes{current}

\logo{\includegraphics[height=0.08\textwidth]{qr.png}}

% 在每个Section前都会加入的Frame
\AtBeginSection[]
{
  \begin{frame}<beamer>
    %\frametitle{Outline}
    \frametitle{教学提纲}
    \setcounter{tocdepth}{3}
    \begin{multicols}{2}
      \tableofcontents[currentsection,currentsubsection]
      %\tableofcontents[currentsection]
    \end{multicols}
  \end{frame}
}
% 在每个Subsection前都会加入的Frame
\AtBeginSubsection[]
{
  \begin{frame}<beamer>
%%\begin{frame}<handout:0>
%% handout:0 表示只在手稿中出现
    \frametitle{教学提纲}
    \setcounter{tocdepth}{3}
    \begin{multicols}{2}
    \tableofcontents[currentsection,currentsubsection]
    \end{multicols}
%% 显示在目录中加亮的当前章节
  \end{frame}
}

% 为当前幻灯片设置背景
%{
%\usebackgroundtemplate{
%\vbox to \paperheight{\vfil\hbox to
%\paperwidth{\hfil\includegraphics[width=2in]{tijmu_charcoal.png}\hfil}\vfil}
%}
\begin{frame}[plain]
  \begin{center}
    {\Huge 疯狂的实验\\}
    \vspace{1cm}
    {\LARGE 天津医科大学\\}
    %\vspace{0.2cm}
    {\LARGE 生物医学工程与技术学院\\}
    \vspace{1cm}
    {\large 2018-2019学年下学期（春）\\ 公共选修课}
  \end{center}
\end{frame}
%}



%\includeonlyframes{current}

\title[生长发育]{第二章\quad 疯狂的实验——生长发育}
\author[Yixf]{伊现富（Yi Xianfu）}
\institute[TIJMU]{天津医科大学（TIJMU）\\ 生物医学工程与技术学院}
\date{2017年3月}

\begin{frame}
  \titlepage
\end{frame}

\begin{frame}[plain,label=current]
  \frametitle{教学提纲}
  \setcounter{tocdepth}{3}
  \begin{multicols}{2}
    \tableofcontents
  \end{multicols}
\end{frame}


\section{母爱}
\begin{frame}
  \frametitle{母爱 | 缘起}
  \begin{block}{当时的观点}
    乳儿实现最佳成长的首要条件是充足的食物和干净的环境。
  \end{block}
  \pause
  \begin{block}{偶然发现}
    将出生不久的小猴与母猴分离：
    \begin{itemize}
      \item 猴子吃住无忧，更壮更重
      \item 总是蜷坐在笼中吮吸手指、呆望远处
      \item 全然不知如何与众多小公/母猴相处
      \item 对铺在笼中的布制衬垫表现出强烈的依恋
    \end{itemize}
  \end{block}
\end{frame}

\begin{frame}
  \frametitle{母爱 | 实验}
  \begin{block}{“母亲机”}
    \begin{itemize}
      \item “绒布妈妈” vs. 挂着奶瓶的“铁丝妈妈”
      \item 各种“怪物妈妈”
    \end{itemize}
  \end{block}
  \pause
  \begin{block}{结果}
    \begin{itemize}
      \item 小猴子们每天紧紧拥抱“绒布妈妈”超过12小时，只有当它们感到饥渴时，才会爬到“铁丝妈妈”身上一会儿
      \item 只要“怪物妈妈”刚一恢复平静，它们就重新跑回紧紧依偎在妈妈身边
    \end{itemize}
  \end{block}
\end{frame}

\begin{frame}
  \frametitle{母爱 | 结论}
  \begin{block}{结论}
    \begin{itemize}
      \item 乳儿对母亲依恋的首要原因是母亲躯体的柔软温暖，与该躯体是否同时作为食物来源并无关系
      \item 身体上的亲近对于孩童的成长非常重要
      \item 幼童追寻母亲并且完全依赖母亲
    \end{itemize}
  \end{block}
  \pause
  \begin{block}{启示}
    \begin{itemize}
      \item 多陪陪孩子（但不要溺爱孩子）
      \item 老人如小孩——多陪陪父母
    \end{itemize}
  \end{block}
\end{frame}

\begin{frame}
  \frametitle{母爱 | 智商}
  \begin{block}{当时的观点}
    \begin{itemize}
      \item 人的智商主要取决于遗传，一生之中基本不变。
      \item 只要满足基本的生理需求，儿童就可以健康成长。
      \item 幼年时期获得过多的温情和关爱是有害的。
    \end{itemize}
  \end{block}
  \vspace{-0.6em}
  \pause
  \begin{block}{“大胆试验”}
    \begin{itemize}
      \item 心理学家哈罗德\textbullet 斯基尔斯
      \item 13名3岁以下的儿童：和智障女性（“妈妈”）待在一起
      \item 12名儿童（比对样本）：留在福利院
    \end{itemize}
  \end{block}
  \vspace{-0.6em}
  \pause
  \begin{block}{结果}
    \begin{itemize}
      \item 儿童期的智商测试：平均提高28分（进步最大的往往有一个固定的“妈妈”） vs. 下降26分
      \item 成年后的生活状态：11人结婚、完整的人生 vs. 9人未婚、1人离异、生活在社会边缘
    \end{itemize}
  \end{block}
\end{frame}

\begin{frame}
  \frametitle{母爱 | 启示}
  \begin{block}{福利院}
    \begin{itemize}
      \item 福利院里有很多看起来发育迟缓、精神涣散的孩子，其实他们并没有先天缺陷，只是他们得到的刺激和关注太少了。
      \item 孤儿福利院原来是滋生智力障碍的温床！
    \end{itemize}
  \end{block}
  \pause
  \begin{block}{智商}
    \begin{itemize}
      \item 关爱和温情让她们走出了迟钝和麻木的世界。
      \item 智力高低并非生来注定，还会受到环境影响，幼年环境尤为重要。
      \item 在智商问题大论战中，“遗传说”和“教育说”仍在激烈交锋。
    \end{itemize}
  \end{block}
\end{frame}

\section{自我意识}
\begin{frame}
  \frametitle{自我意识 | 实验}
  \begin{block}{问题}
    动物是否具有自我意识？
  \end{block}
  \pause
  \begin{block}{实验}
    \begin{itemize}
      \item 实验组：照镜子——涂红点——照镜子
      \item 对照组：涂红点——照镜子
    \end{itemize}
  \end{block}
  \pause
  \begin{block}{引申}
    许多动物行为学家在自己孩子身上做的第一项科学实验
  \end{block}
\end{frame}

\begin{frame}
  \frametitle{自我意识 | 启示}
  \begin{block}{本质}
    \begin{itemize}
      \item 动物拥有辨认镜中自我的能力，到底意味着什么？
      \item 是否意味着它已经可以明确认识到那是自己？
      \item 是否意味着它也能同样辨别出同类？
      \item 是否意味着它也会像人类那般做出计划或者撒谎？
    \end{itemize}
  \end{block}
  \pause
  \begin{block}{启示}
    \begin{itemize}
      \item 在科学史上偶然间冒出绝妙灵感的事例层出不穷
      \item 许多科学实验能够带给人有趣的答案，然而却未必真正触及了问题的本质
    \end{itemize}
  \end{block}
\end{frame}

\section{幼儿喂养}
\begin{frame}
  \frametitle{幼儿喂养 | 问题}
  \begin{block}{旧式的育儿观念}
    从母乳改为成人食品的调整过程需要持续3-4年。
  \end{block}
  \pause
  \begin{block}{营养学}
    动物（包括人）能够在丰富的食物供应中凭借直觉选择出最适合自身成长发展的食物，还是必须总从生物化学家制定的食谱，先由他们分析每种食物的营养成分，让后再做选择？
  \end{block}
\end{frame}

\begin{frame}
  \frametitle{幼儿喂养 | 实验}
  \begin{block}{实验}
    \begin{itemize}
      \item 实验对象：15名6个月到4岁半不等的婴幼儿
      \item 实验食品：苹果、菠萝泥、西红柿、煎土豆、煮熟的小麦、玉米、燕麦和黑麦、做好并剁碎的牛肉、骨髓、动物的脑、肝、肾、剁碎的鱼肉、鸡蛋、食盐、水、不同种类的牛奶和橙汁，……
      \item 实验控制：不接受父母的干预和儿科医生的指令
      \item 实验测量：称量吃掉的分量，检查生长发育状况
    \end{itemize}
  \end{block}
\end{frame}

\begin{frame}
  \frametitle{幼儿喂养 | 结果}
  \begin{block}{结果}
    \begin{itemize}
      \item 孩子们的生长发育完全正常，没有任何营养缺乏的迹象。
      \item 每个儿童的首选“菜单”不仅差异明显，还会固定偏爱某种食品。
      \item 孩子们吃水果、肉类、蛋类、脂肪类食物的数量远远超过了当时儿科医生的建议量，吃粮食和蔬菜的数量则少于建议量。
      \item 菠菜几乎遭到所有孩子的厌弃，圆白菜和生菜也不太受欢迎。
      \item 乍一看来，孩子们的营养搭配乱七八糟；仔细观察，他们都做出了明知合理的营养选择：蛋白质、脂肪和碳水化合物的数量都在标准范围之内。
    \end{itemize}
  \end{block}
  \pause
  \begin{block}{结论}
    \begin{itemize}
      \item 所有人都遵循的标准食谱未必就是最优的膳食搭配。
      \item 儿童可以毫无障碍地消化成人食品并正常地长大。
    \end{itemize}
  \end{block}
\end{frame}

\section{思维发展}
\begin{frame}
  \frametitle{思维发展 | 现象}
  \begin{block}{现象}
    \begin{itemize}
      \item 数量守恒——某个东西被分成若干部分或者改变了形状，其总量并不会突然增多或者减少
      \item 客体永恒性——事物具有恒定不变的特征
      \item 空间参考体系——反映出空间观念的发展和对空间能力的运用
    \end{itemize}
  \end{block}
  \pause
  \begin{block}{实验}
    \begin{itemize}
      \item 让\textbullet 皮亚杰，不同年龄段的儿童
      \item 数量守恒实验
      \item 水平面任务
    \end{itemize}
  \end{block}
\end{frame}

\begin{frame}
  \frametitle{思维发展 | 启示}
  \begin{block}{儿童思维发展理论}
    \begin{itemize}
      \item 儿童要在特定的年龄才能到到某个阶段，从儿童所犯的典型错误能够识别出他们正处于什么阶段。
      \item 在前运算阶段（2-7岁），儿童的判断主要由感知决定。
      \item 在具体运算阶段（7-12）岁，儿童开始按照逻辑规则去思考。
    \end{itemize}
  \end{block}
  \pause
  \begin{block}{启示}
    \begin{itemize}
      \item 儿童的思维错误就是宝贵的知识源泉。
      \item 皮亚杰的实验绝对堪称心理学界最重要和最富创造力的实验。
      \item 皮亚杰虽然是极为杰出的思想者，却不是特别严谨的实验员。
      \item 皮亚杰通过水平面任务测试说明了儿童空间想象力的发展。
    \end{itemize}
  \end{block}
\end{frame}

\begin{frame}
  \frametitle{思维发展 | 水平面任务}
  \begin{block}{儿童}
    \begin{itemize}
      \item 5岁儿童画出的水平面总是与杯壁垂直。
      \item 6或7岁的儿童发现，垂直是不对的，但水平面还是被画成了倾斜状。
      \item 大约9岁的儿童才能发现正确的答案：水平面总是呈现水平状态，也就是与桌面平行。
    \end{itemize}
  \end{block}
  \pause
  \begin{block}{成人}
    \begin{itemize}
      \item 女性在水平面任务中的表现明显比男性差。
      \item 错误率随着经验的增加而提高？（经验诱导人们将玻璃杯作为参照系）
      \item 会写中文的人在水平面任务中得分更高。
    \end{itemize}
  \end{block}
\end{frame}

\begin{frame}
  \frametitle{思维发展 | 幼儿思维能力}
  \begin{block}{问题}
    人类的哪些技能是生来具备的，哪些技能是后天习得的，先天基础和外部环境，哪个起着决定作用？
  \end{block}
  \pause
  \begin{block}{当时的观点}
    新生儿就像一张干净的白纸，世界在他们眼中只是一些或明或暗的色块，随着“观看”经验与日俱增，他们才逐渐学会分辨不同的事物。
  \end{block}
  \pause
  \begin{block}{实验}
    罗伯特\textbullet 法恩茨，22个婴儿（1-15周不等），观看各种形状，观察条纹板和纯色板，观察加法题
  \end{block}
\end{frame}

\begin{frame}
  \frametitle{思维发展 | 幼儿思维能力}
  \begin{block}{结果}
    \begin{itemize}
      \item 都喜欢观察比较复杂的图形——婴儿天生就具备区分这些图形的能力。
      \item 一个月大的婴儿能够识别3mm宽的条纹，半岁大的可以识别再细10倍的条纹。
      \item 看球体的时间比看圆环的时间更长。
      \item 人脸更受欢迎。
      \item 婴儿能够运算。
    \end{itemize}
  \end{block}
\end{frame}

\section{延迟满足}
\begin{frame}
  \frametitle{延迟满足 | 实验}
  \begin{block}{预测}
    如何预测一名4岁儿童的未来——他能否成为一个健全而知足的人（优异成绩、很多朋友、吸毒、和谐伴侣……）？
  \end{block}
  \pause
  \begin{block}{棉花糖测试}
    让他选择是现在吃到1块棉花糖，还是以后吃到2块。——孩子愿意为2块棉花糖等待的时间越长，他在以后的人生中就表现得越好。
  \end{block}
\end{frame}

\begin{frame}
  \frametitle{延迟满足 | 启示}
  \begin{block}{延迟满足}
    应该马上服从自己的需求，还是为了一个更高的目标推迟这些需求呢？
  \end{block}
  \vspace{-0.5em}
  \pause
  \begin{block}{结果}
    \begin{itemize}
      \item 成长过程中没有父亲陪伴的儿童往往不想等待更大的奖励。
      \item 在棉花糖测试里显示出耐心的儿童，在学校的表现明显好于其他人，也较少出现各种问题。
      \item 4-6岁完成的棉花糖测试以出人意料的准确度预测了被试儿童10年后的许多特征。
    \end{itemize}
  \end{block}
  \vspace{-0.5em}
  \pause
  \begin{block}{启示}
    \begin{itemize}
      \item 自发地延迟满足是迈向成熟的重要一步。
      \item “为长期目标放弃短期诱惑”是人生中最应该掌握的重要能力之一。
      \item 实验结果深受实验组织方式和奖励类型的影响，而且归根结底也只反映出一种统计学上的倾向，就个体案例来说指导作用有限。
    \end{itemize}
  \end{block}
\end{frame}

\section{史海撷华}
\begin{frame}
  \frametitle{史海撷华 | 通往香蕉的阶梯}
  \begin{block}{人}
    “把第二个建筑材料放到第一个上只是对把第一个建筑材料放到地上这一行为的重复”。
  \end{block}
  \pause
  \begin{block}{黑猩猩}
    黑猩猩需要解决的问题“分为两个截然不同的部分：一部分它可以轻易解决，而另一部分解决时却遇到了相当大的困难”。容易解决的部分是把一只箱子推到香蕉下面，困难的部分是把第二只箱子叠垛到先前的那只上面。
  \end{block}
  \pause
  \begin{block}{启示}
    人和黑猩猩相似的行为并不能证明它们相似的思维方式。
  \end{block}
\end{frame}

\begin{frame}
  \frametitle{史海撷华 | 猴妹妹}
  \begin{block}{经验/现象}
    \begin{itemize}
      \item 从小看大，三岁看老
      \item 三岁定八十
      \item 狼/熊/虎/猴孩
    \end{itemize}
  \end{block}
  \pause
  \begin{block}{假说/猜想}
    \begin{itemize}
      \item 想要把生命早期已形成的习性再次革除并非易事？
      \item 在成长过程中，是本性还是文化，是环境还是遗传，起了决定性作用？
    \end{itemize}
  \end{block}
\end{frame}

\begin{frame}
  \frametitle{史海撷华 | 猴妹妹}
  \begin{block}{实验}
    \begin{itemize}
      \item 10个月大的唐纳德，7个月大的古亚
      \item 猴子和人类的孩子被用完全相同的方式对待9个月
    \end{itemize}
  \end{block}
  \pause
  \begin{block}{结果/后果}
    \begin{itemize}
      \item 古亚表现出对人类环境惊人的适应性；唐纳德是一个更优秀的模仿者。
      \item 本想把一只猴子教育成人，却把一个人教成了猴子。
      \item 唐纳德，一名精神病医生，（45年）自杀
    \end{itemize}
  \end{block}
\end{frame}

\begin{frame}
  \frametitle{史海撷华 | 先天 vs. 教育}
  \begin{block}{实验人员}
    \begin{itemize}
      \item 实验者：莫特尔\textbullet 麦格劳（发现“潜游反射”）
      \item 实验对象：约翰尼、吉米（同卵双胞胎？异卵双胞胎！）
    \end{itemize}
  \end{block}
  \pause
  \begin{block}{问题}
    \begin{itemize}
      \item “促进措施”对儿童运动技能发展的影响
      \item 先天因素和教育因素哪个作用更大
    \end{itemize}
  \end{block}
  \pause
  \begin{block}{成果：约翰尼}
    \begin{itemize}
      \item 13个月，会溜旱冰
      \item 15个月，完成从1.5米高的跳板上俯身跃下的动作
      \item 17个月，在水下潜游4米
      \item 21个月，能从1.6米高的台子上爬下来
      \item 22个月，毫不费力地攀爬上70度的陡峭斜坡
    \end{itemize}
  \end{block}
\end{frame}

\begin{frame}
  \frametitle{史海撷华 | 玩具偏好}
  \begin{block}{现象}
    儿童做事通常只有3分钟热度，可是，不同性别的儿童对特定玩具的偏好却稳定而持久。
  \end{block}
  \vspace{-0.5em}
  \pause
  \begin{block}{实验}
    \begin{itemize}
      \item 88只绿长尾猴——44只雌性和44只雄性
      \item 6中玩具：2种典型的男性玩具（球和警车），2种典型的女性玩具（布娃娃和厨房用锅），2种中性玩具（图画书和毛绒玩具狗）
      \item 结果：雄性猴子玩球和警车的时间是雌性的2倍；雌性猴子玩布娃娃和厨房用锅的时间是雄性的2倍；图画书和毛绒玩具狗在雌性和雄性猴子中的受欢迎程度基本相当。
    \end{itemize}
  \end{block}
  \vspace{-0.5em}
  \pause
  \begin{block}{结论}
    不同性别的个体对不同玩具的偏好并不仅仅由家长和电视广告决定，生物学也在发挥作用。
  \end{block}
\end{frame}



\section*{Acknowledgements}
\begin{frame}
  \frametitle{Powered by}
  \begin{center}
    \includegraphics[width=9cm]{power.png}
  \end{center}
\end{frame}

\end{document}

